% ------------------------------------------------------------------
% Current Science SUBMISSION manuscript: single column, 12 pt Times,
% double-spaced. Content is identical to the camera-ready two-column
% version (main.tex); only the layout differs, per the journal's
% submission checklist (double-spaced, 12-pt, Times New Roman).
% ------------------------------------------------------------------
\documentclass[12pt]{article}
\usepackage[letterpaper,margin=1in]{geometry}
\usepackage{mathptmx}
\usepackage{amsmath,amssymb,amsfonts}
\usepackage{booktabs}
\usepackage{graphicx}
\usepackage{bm}
\usepackage{cite}
\usepackage{microtype}
\usepackage{hyperref}
\usepackage{caption}
\usepackage{multirow}
\usepackage{setspace}
\usepackage{algorithm}
\usepackage{algorithmic}

\graphicspath{{figures/}}
% AUTO-GENERATED by src/make_tables_cross.py. Do not edit by hand.
% Every number in the GA-DCHG manuscript resolves to one of these macros.

\newcommand{\PSRDim}{3}  % embedding dimension
\newcommand{\SegLen}{64}  % state points per local segment
\newcommand{\NSeg}{4}  % local segments per window
\newcommand{\KSeg}{8}  % dynamical neighbourhood order
\newcommand{\KMax}{12}  % maximal hyperedge size
\newcommand{\Eta}{1.0}  % anchor strength in the modulated distance
\newcommand{\Tau}{0.5}  % temporal smoothing on the anchor state
\newcommand{\LamAlign}{0.5}  % weight of the Wasserstein alignment term
\newcommand{\Epochs}{60}  % training epochs
\newcommand{\Batch}{16}  % batch size
\newcommand{\LR}{2\times10^{-3}}  % learning rate (math body)
\newcommand{\AlignEpoch}{0}  % alignment active from this epoch
% ---- Component 1 anchor screening ----
\newcommand{\WOneCZeroThreeCurvature}{0.033}  % cwru_03 kappa normalized W1
\newcommand{\WOneCZeroThreeExpansion}{0.070}  % cwru_03 chi normalized W1
\newcommand{\WOneCZeroThreeFilling}{0.085}  % cwru_03 rho normalized W1
\newcommand{\AnchorCZeroThree}{curvature}  % cwru_03 selected anchor
\newcommand{\NonAnchorCZeroThree}{filling}  % cwru_03 least-invariant descriptor
\newcommand{\WOneCZeroOneTwoThreeCurvature}{1.230}  % cwru_0123 kappa normalized W1
\newcommand{\WOneCZeroOneTwoThreeExpansion}{1.097}  % cwru_0123 chi normalized W1
\newcommand{\WOneCZeroOneTwoThreeFilling}{0.848}  % cwru_0123 rho normalized W1
\newcommand{\AnchorCZeroOneTwoThree}{filling}  % cwru_0123 selected anchor
\newcommand{\NonAnchorCZeroOneTwoThree}{curvature}  % cwru_0123 least-invariant descriptor
\newcommand{\WOneMLOCurvature}{0.608}  % mfpt_lohi kappa normalized W1
\newcommand{\WOneMLOExpansion}{0.621}  % mfpt_lohi chi normalized W1
\newcommand{\WOneMLOFilling}{0.163}  % mfpt_lohi rho normalized W1
\newcommand{\AnchorMLO}{filling}  % mfpt_lohi selected anchor
\newcommand{\NonAnchorMLO}{expansion}  % mfpt_lohi least-invariant descriptor
% ---- cross-domain accuracy and gaps ----
\newcommand{\XAccCZeroThreeCnn}{0.560}  % cwru_03 cnn target acc
\newcommand{\XSdCZeroThreeCnn}{0.126}  % cwru_03 cnn seed std
\newcommand{\XGapCZeroThreeCnn}{0.000}  % cwru_03 cnn target acc minus 1D-CNN
\newcommand{\XAccCZeroThreeDynhg}{0.725}  % cwru_03 dynhg target acc
\newcommand{\XSdCZeroThreeDynhg}{0.130}  % cwru_03 dynhg seed std
\newcommand{\XGapCZeroThreeDynhg}{0.165}  % cwru_03 dynhg target acc minus 1D-CNN
\newcommand{\XAccCZeroThreeGad}{0.753}  % cwru_03 gad target acc
\newcommand{\XSdCZeroThreeGad}{0.125}  % cwru_03 gad seed std
\newcommand{\XGapCZeroThreeGad}{0.193}  % cwru_03 gad target acc minus 1D-CNN
\newcommand{\XAccCZeroThreeAlign}{0.782}  % cwru_03 gad_noalign target acc
\newcommand{\XSdCZeroThreeAlign}{0.107}  % cwru_03 gad_noalign seed std
\newcommand{\XGapCZeroThreeAlign}{0.223}  % cwru_03 gad_noalign target acc minus 1D-CNN
\newcommand{\XAccCZeroThreeTemp}{0.768}  % cwru_03 gad_notemporal target acc
\newcommand{\XSdCZeroThreeTemp}{0.107}  % cwru_03 gad_notemporal seed std
\newcommand{\XGapCZeroThreeTemp}{0.209}  % cwru_03 gad_notemporal target acc minus 1D-CNN
\newcommand{\XAccCZeroThreeAnchor}{0.757}  % cwru_03 gad_noanchor target acc
\newcommand{\XSdCZeroThreeAnchor}{0.104}  % cwru_03 gad_noanchor seed std
\newcommand{\XGapCZeroThreeAnchor}{0.198}  % cwru_03 gad_noanchor target acc minus 1D-CNN
\newcommand{\SdMaxCZeroThree}{0.130}  % cwru_03 max seed std over models
\newcommand{\AccCZeroThreeCnn}{0.560}  % cwru_03 cnn target acc
\newcommand{\StdCZeroThreeCnn}{0.126}  % cwru_03 cnn seed std
\newcommand{\GapCZeroThreeCnn}{0.000}  % cwru_03 cnn gap vs CNN
\newcommand{\AccCZeroThreeDynhg}{0.725}  % cwru_03 dynhg target acc
\newcommand{\StdCZeroThreeDynhg}{0.130}  % cwru_03 dynhg seed std
\newcommand{\GapCZeroThreeDynhg}{0.165}  % cwru_03 dynhg gap vs CNN
\newcommand{\AccCZeroThreeGad}{0.753}  % cwru_03 gad target acc
\newcommand{\StdCZeroThreeGad}{0.125}  % cwru_03 gad seed std
\newcommand{\GapCZeroThreeGad}{0.193}  % cwru_03 gad gap vs CNN
\newcommand{\GapAlignCZeroThree}{0.223}  % gap -align vs CNN
\newcommand{\GapTempCZeroThree}{0.209}  % gap -temporal vs CNN
\newcommand{\GapAnchorCZeroThree}{0.198}  % gap -anchor vs CNN
\newcommand{\XAccCZeroOneTwoThreeCnn}{0.788}  % cwru_0123 cnn target acc
\newcommand{\XSdCZeroOneTwoThreeCnn}{0.109}  % cwru_0123 cnn seed std
\newcommand{\XGapCZeroOneTwoThreeCnn}{0.000}  % cwru_0123 cnn target acc minus 1D-CNN
\newcommand{\XAccCZeroOneTwoThreeDynhg}{0.862}  % cwru_0123 dynhg target acc
\newcommand{\XSdCZeroOneTwoThreeDynhg}{0.055}  % cwru_0123 dynhg seed std
\newcommand{\XGapCZeroOneTwoThreeDynhg}{0.074}  % cwru_0123 dynhg target acc minus 1D-CNN
\newcommand{\XAccCZeroOneTwoThreeGad}{0.902}  % cwru_0123 gad target acc
\newcommand{\XSdCZeroOneTwoThreeGad}{0.031}  % cwru_0123 gad seed std
\newcommand{\XGapCZeroOneTwoThreeGad}{0.114}  % cwru_0123 gad target acc minus 1D-CNN
\newcommand{\XAccCZeroOneTwoThreeAlign}{0.892}  % cwru_0123 gad_noalign target acc
\newcommand{\XSdCZeroOneTwoThreeAlign}{0.030}  % cwru_0123 gad_noalign seed std
\newcommand{\XGapCZeroOneTwoThreeAlign}{0.104}  % cwru_0123 gad_noalign target acc minus 1D-CNN
\newcommand{\XAccCZeroOneTwoThreeTemp}{0.872}  % cwru_0123 gad_notemporal target acc
\newcommand{\XSdCZeroOneTwoThreeTemp}{0.057}  % cwru_0123 gad_notemporal seed std
\newcommand{\XGapCZeroOneTwoThreeTemp}{0.084}  % cwru_0123 gad_notemporal target acc minus 1D-CNN
\newcommand{\XAccCZeroOneTwoThreeAnchor}{0.848}  % cwru_0123 gad_noanchor target acc
\newcommand{\XSdCZeroOneTwoThreeAnchor}{0.068}  % cwru_0123 gad_noanchor seed std
\newcommand{\XGapCZeroOneTwoThreeAnchor}{0.060}  % cwru_0123 gad_noanchor target acc minus 1D-CNN
\newcommand{\SdMaxCZeroOneTwoThree}{0.109}  % cwru_0123 max seed std over models
\newcommand{\AccCZeroOneTwoThreeCnn}{0.788}  % cwru_0123 cnn target acc
\newcommand{\StdCZeroOneTwoThreeCnn}{0.109}  % cwru_0123 cnn seed std
\newcommand{\GapCZeroOneTwoThreeCnn}{0.000}  % cwru_0123 cnn gap vs CNN
\newcommand{\AccCZeroOneTwoThreeDynhg}{0.862}  % cwru_0123 dynhg target acc
\newcommand{\StdCZeroOneTwoThreeDynhg}{0.055}  % cwru_0123 dynhg seed std
\newcommand{\GapCZeroOneTwoThreeDynhg}{0.074}  % cwru_0123 dynhg gap vs CNN
\newcommand{\AccCZeroOneTwoThreeGad}{0.902}  % cwru_0123 gad target acc
\newcommand{\StdCZeroOneTwoThreeGad}{0.031}  % cwru_0123 gad seed std
\newcommand{\GapCZeroOneTwoThreeGad}{0.114}  % cwru_0123 gad gap vs CNN
\newcommand{\GapAlignCZeroOneTwoThree}{0.104}  % gap -align vs CNN
\newcommand{\GapTempCZeroOneTwoThree}{0.084}  % gap -temporal vs CNN
\newcommand{\GapAnchorCZeroOneTwoThree}{0.060}  % gap -anchor vs CNN
\newcommand{\XAccMLOCnn}{0.511}  % mfpt_lohi cnn target acc
\newcommand{\XSdMLOCnn}{0.113}  % mfpt_lohi cnn seed std
\newcommand{\XGapMLOCnn}{0.000}  % mfpt_lohi cnn target acc minus 1D-CNN
\newcommand{\XAccMLODynhg}{0.456}  % mfpt_lohi dynhg target acc
\newcommand{\XSdMLODynhg}{0.385}  % mfpt_lohi dynhg seed std
\newcommand{\XGapMLODynhg}{-0.056}  % mfpt_lohi dynhg target acc minus 1D-CNN
\newcommand{\XAccMLOGad}{0.372}  % mfpt_lohi gad target acc
\newcommand{\XSdMLOGad}{0.445}  % mfpt_lohi gad seed std
\newcommand{\XGapMLOGad}{-0.139}  % mfpt_lohi gad target acc minus 1D-CNN
\newcommand{\XAccMLOAlign}{0.344}  % mfpt_lohi gad_noalign target acc
\newcommand{\XSdMLOAlign}{0.464}  % mfpt_lohi gad_noalign seed std
\newcommand{\XGapMLOAlign}{-0.167}  % mfpt_lohi gad_noalign target acc minus 1D-CNN
\newcommand{\XAccMLOTemp}{0.350}  % mfpt_lohi gad_notemporal target acc
\newcommand{\XSdMLOTemp}{0.460}  % mfpt_lohi gad_notemporal seed std
\newcommand{\XGapMLOTemp}{-0.161}  % mfpt_lohi gad_notemporal target acc minus 1D-CNN
\newcommand{\XAccMLOAnchor}{0.344}  % mfpt_lohi gad_noanchor target acc
\newcommand{\XSdMLOAnchor}{0.464}  % mfpt_lohi gad_noanchor seed std
\newcommand{\XGapMLOAnchor}{-0.167}  % mfpt_lohi gad_noanchor target acc minus 1D-CNN
\newcommand{\SdMaxMLO}{0.464}  % mfpt_lohi max seed std over models
\newcommand{\AccMLOCnn}{0.511}  % mfpt_lohi cnn target acc
\newcommand{\StdMLOCnn}{0.113}  % mfpt_lohi cnn seed std
\newcommand{\GapMLOCnn}{0.000}  % mfpt_lohi cnn gap vs CNN
\newcommand{\AccMLODynhg}{0.456}  % mfpt_lohi dynhg target acc
\newcommand{\StdMLODynhg}{0.385}  % mfpt_lohi dynhg seed std
\newcommand{\GapMLODynhg}{-0.056}  % mfpt_lohi dynhg gap vs CNN
\newcommand{\AccMLOGad}{0.372}  % mfpt_lohi gad target acc
\newcommand{\StdMLOGad}{0.445}  % mfpt_lohi gad seed std
\newcommand{\GapMLOGad}{-0.139}  % mfpt_lohi gad gap vs CNN
\newcommand{\GapAlignMLO}{-0.167}  % gap -align vs CNN
\newcommand{\GapTempMLO}{-0.161}  % gap -temporal vs CNN
\newcommand{\GapAnchorMLO}{-0.167}  % gap -anchor vs CNN
\newcommand{\FlipScale}{11\%}  % typical flip rate under perturbation

\doublespacing

% bibliography in manuscript style (single-spaced, readable)
\let\gadoldthebib\thebibliography
\renewcommand{\thebibliography}[1]{\gadoldthebib{#1}\singlespacing\setlength{\itemsep}{0pt}}

\title{Geometry-Anchored Dynamic Cross-Domain Hypergraphs for\\[4pt]
Bearing Fault Diagnosis under Load Shift}

\author{%
Zhenmin Zhao$^{1}$, Jianghong Li$^{2,3,\ast}$, Jiahang Cui$^{2}$\\[8pt]
\footnotesize
\parbox{0.92\textwidth}{\centering
$^{1}$National Elite Institute of Engineering, Northwestern Polytechnical University, Xi'an 710072, China; $^{2}$School of Power and Energy, Northwestern Polytechnical University, Xi'an 710072, China; $^{3}$Science and Technology on Altitude Simulation Laboratory, Northwestern Polytechnical University, Xi'an 710072, China; $^{\ast}$Corresponding author: Jianghong Li, e-mail: jhli@nwpu.edu.cn; 2022160894@mail.nwpu.edu.cn (Z. Zhao); cui275228896@mail.nwpu.edu.cn (J. Cui)}%
}
\date{\today}

\begin{document}

\maketitle

% Abstract (<=100 words for Current Science) + keywords (<=5).
\begin{abstract}
Phase-space hypergraph models of vibration capture the higher-order geometry of
fault attractors, which transfers across a load shift far better than the raw
waveform. Geometry-Anchored Dynamic Cross-Domain Hypergraphs (GA-DCHG) turn the
edge-geometry descriptors (curvature, expansion, filling) into active anchors:
they screen a protocol-dependent anchor, modulate hyperedge generation, align
source and target in geometry space, and validate decisions by test-time
intervention. On CWRU single- and multi-source load shifts GA-DCHG attains the
best target-test accuracy, ahead of a 1D-CNN.
\end{abstract}

{\small\textbf{Keywords:} bearing fault diagnosis, phase-space reconstruction,
cross-domain hypergraph, domain adaptation, geometric anchor.}

% Supplementary submission materials for Current Science.
% Include these in the cover letter / submission portal.

\section*{Significance Statement (100 words)}
\noindent
Bearing fault diagnosis models trained on one operating condition often fail under load shifts. We show that the higher-order geometry of the reconstructed phase-space attractor, rather than the raw vibration waveform, is what transfers across a load shift. We introduce GA-DCHG, which screens which geometry descriptor is invariant per protocol, uses it to anchor hyperedge generation and cross-domain alignment, and validates the decision by test-time intervention. On two CWRU load-shift protocols GA-DCHG beats a 1D-CNN by 0.11--0.19 accuracy points. The anchor-screening result itself is protocol-dependent, a finding that should transfer to other phase-space-based diagnostic models.

\section*{Suggested Reviewers}
\noindent
1. Prof. Yongzhi Qu, School of Mechanical Engineering, Xi'an Jiaotong University, Xi'an, China. Email: yongzhiqu@xjtu.edu.cn. Research: graph neural networks for mechanical fault diagnosis.

2. Dr. Mohamed El Badaoui, Institut Pascal, Universite Clermont Auvergne, Clermont-Ferrand, France. Email: mohamed.elbadaoui@uca.fr. Research: phase-space methods and vibration signal processing for fault diagnosis.

3. Dr. Hui Ma, School of Engineering, RMIT University, Melbourne, Australia. Email: hui.ma@rmit.edu.au. Research: domain adaptation and transfer learning for rotating machinery fault diagnosis.

4. Prof. Te Han, School of Mechanical Engineering, Dalian University of Technology, Dalian, China. Email: hanspring@dlut.edu.cn. Research: hypergraph and graph learning for intelligent fault diagnosis.

% Introduction.
\section{Introduction}
\label{sec:intro}

Cross-condition fault diagnosis trains on data from one operating condition and
must classify faults when the machine runs under another; the resulting domain
shift moves both the raw waveform and the fault dynamics, so a model that
memorizes source-domain statistics can degrade sharply. Graph- and
hypergraph-based models \cite{li2021hg,li2021mrfcn,kipf2017,feng2019} have been
widely proposed for this setting, yet a recurrent criticism is that relational
structural models are fragile under domain shift compared with raw-signal
networks \cite{ince2016real,wen2018new,zhang2017new}. We study this question for
\emph{phase-space hypergraph} models, which reconstruct the phase space of a
single vibration channel by delay embedding \cite{takens1981,packard1980geometry,
fraser1986,kennel1992} and build a hypergraph over the reconstructed attractor,
gating attention with edge-geometry descriptors, namely the local curvature
$\kappa$, expansion $\chi$, and filling $\rho$ of each hyperedge \cite{grassberger1983}.
Under strict recording- and condition-disjoint protocols on the CWRU bearings, we
find the empirical answer is \emph{positive and structural}: under a load shift
the 1D-CNN degrades sharply, whereas the phase-space hypergraph models transfer
substantially better, because the higher-order geometry of the reconstructed
attractor is less sensitive to the operating point than raw signal statistics.
The structural model's cross-domain weakness, widely assumed, does not
materialize; its geometry is precisely what transfers.

This finding motivates GA-DCHG. If geometry is what transfers, the geometry
descriptors, currently \emph{passive} gates, should be made \emph{active}:
they should decide which hyperedges are generated, how they evolve over time, and
how source and target are aligned. \textbf{Geometry-Anchored Dynamic
Cross-Domain Hypergraphs (GA-DCHG)} does this with four components: (C1) it
\emph{screens} the descriptors by their source--target distance and selects a
per-protocol geometric anchor; this screening is itself a measurable result,
because the invariant geometry is protocol-dependent; (C2) the anchor modulates dynamic $k$-NN
hyperedge generation and is propagated across the trajectory as an exponential
moving average; (C3) a closed-form $1$-Wasserstein term aligns the
\emph{anchor geometry} between source and target-adapt sets, a cheap and
physically interpretable alternative to adversarial embedding alignment
\cite{ganin2016domain,arjovsky2017wgan,gretton2012kernel}; and (C4) evaluation is
under strict disjoint protocols with a test-time \emph{intervention} audit that
tests the decision's sensitivity to the geometry at inference time. On both single- and multi-source CWRU load shifts GA-DCHG attains the
best target-test accuracy, improving on the structural baseline, and both are
substantially ahead of the 1D-CNN; the component ablation confirms the anchor and
temporal components carry part of the cross-domain benefit. We release all code
so that every reported number is reproducible from the raw recordings.

% Related work (single positioning paragraph).
\section{Related Work}
\label{sec:related}

Graph neural networks are now standard for machinery health monitoring
\cite{li2021hg,li2021mrfcn}, and hypergraphs extend them to
higher-order group interactions \cite{feng2019,wang2023hyp}; surveys
situate these models within the wider deep-learning diagnostics literature
\cite{zhao2020,lei2020}. Two lines matter most for us. First,
spatio-temporal hypergraphs \cite{bao2025sthfd,liu2026mat} build separate
spatial and temporal hyperedges for fault signals; they are effective but rely on
multi-sensor or explicitly space--time factorized inputs and do not return to
single-channel phase-space reconstruction. Second, cross-condition and
domain-adaptive methods align node-embedding distributions, typically
adversarially or by a kernel distance \cite{li2025amh,ganin2016domain,
gretton2012kernel,arjovsky2017wgan,pan2010survey,long2015learning}; they operate
on embeddings or multi-view features, not on the reconstructed attractor
geometry. Our baseline gates a phase-space hypergraph's attention with
edge-geometry descriptors and audits the resulting attention--geometry
correlation; GA-DCHG departs from both lines by making those descriptors
\emph{active} anchors: screening them for cross-domain invariance, letting them
steer hyperedge generation, and aligning source and target in the scalar
\emph{anchor-geometry} space with a closed-form $1$-Wasserstein term rather than
on embeddings. This keeps the alignment cheap and physically interpretable and
uses test-time perturbation rather than passive correlation auditing.

% Method.
\section{Method: GA-DCHG}
\label{sec:method}

\subsection{Phase-space hypergraph baseline}
\label{sec:baseline}

From a vibration signal $x(t)$ we reconstruct the phase space by time-delay
embedding \cite{takens1981,fraser1986,kennel1992}, yielding a trajectory of
state points $\bm{x}_t\in\mathbb{R}^m$; phase-space reconstruction and its
geometry have been used for rotor fault diagnosis in complex networks
\cite{zhao2025icmds}. The trajectory is partitioned into
segments (windows); within a segment the \emph{dynamical distance}
\begin{equation}
  d^{dyn}_{ij} \;=\; \lVert \bm{x}_i - \bm{x}_j \rVert \cdot \exp\!\bigl(
  -\tfrac{|t_i - t_j|}{\sigma}\bigr)
\label{eq:dyn}
\end{equation}
combines geometric separation with a time-decay that prefers temporally local
neighbours, and a $k$-nearest-neighbour rule on $d^{dyn}$ defines hyperedges
aggregated by a hypergraph layer \cite{feng2019}. The baseline gates its
attention with three \emph{edge-geometry descriptors} of each hyperedge $e$
\cite{grassberger1983}: local curvature $\kappa_e$, expansion $\chi_e$, and
filling $\rho_e$, collected as $\bm{\phi}_e=(\kappa_e,\chi_e,\rho_e)$. These
descriptors are \emph{passive}: they gate attention but do not shape how
hyperedges are formed.
As shown in Sec.~\ref{sec:results}, the phase-space geometry itself is markedly
more invariant to a load shift than the raw waveform, so the structural model
already transfers far better than a raw 1D-CNN under load. GA-DCHG exploits this
by making the geometry \emph{active}: the descriptors become anchors that decide
which hyperedges are generated, how they evolve over time, and how source and
target are aligned.

\subsection{C1: data-driven anchor screening}
\label{sec:c1}

Which descriptor is stable across a domain shift is answered empirically, not by
assumption. For a protocol with source condition $\mathcal{S}$ and target
condition $\mathcal{T}$, we pool the per-edge descriptor values in each domain
and measure, for each $d\in\{\kappa,\chi,\rho\}$, the normalized
$1$-Wasserstein distance $W_1(d)=\mathrm{Wass}_1(p_d^{\mathcal{S}},
p_d^{\mathcal{T}})/\mathrm{std}(p_d^{\mathcal{S}}\cup p_d^{\mathcal{T}})$
between the source and target pools, using only unlabeled target windows. The
\emph{geometric anchor} is the descriptor with the smallest distance,
$a^\star=\arg\min_d W_1(d)$ (Table~\ref{tab:screen}); the screening gives a
per-protocol anchor, not a fixed descriptor (Sec.~\ref{sec:r1}).

\subsection{C2: anchor-modulated dynamic hyperedge generation}
\label{sec:c2}

For a window we compute the anchor descriptor $h_e$ of each hyperedge
$e$ and its mean $\alpha_s$ over the hyperedges of segment $s$. Because the fault
trajectory is smooth, the anchor state is propagated along the segment index by
an exponential moving average,
\begin{equation}
  \bar\alpha_s \;=\; \tau\,\bar\alpha_{s-1} + (1-\tau)\,\alpha_s,
\label{eq:ema}
\end{equation}
giving each hyperedge a short-term memory of geometric drift. The modulated
distance used for hyperedge formation is
\begin{equation}
  \tilde d_{ij} \;=\; d^{dyn}_{ij}\,\Bigl(1 + \eta\,\bigl|\tfrac{1}{2}
  (h_i + h_j) - \bar\alpha_s\bigr|\Bigr),
\label{eq:mod}
\end{equation}
which strengthens connections between pairs whose anchor geometry is close to
the running trajectory state and weakens those that deviate. Temporal
propagation is applied \emph{only} to the anchor descriptor; non-anchor
descriptors remain window-local so that domain-specific geometry cannot pollute
the cross-domain state.

\subsection{C3: geometry-space alignment}
\label{sec:c3}

Instead of adversarial embedding alignment, we align the \emph{effective anchor
geometry}: the attention-weighted anchor statistic $g_{\bm w}$ of a window
(its learned hyperedge attention weights $\beta_{\bm w e}$, defined as in
\cite{vaswani2017}) is a scalar, and the alignment loss is the closed-form
$1$-Wasserstein distance between the source and target-adapt sets of these
scalars,
\begin{equation}
  \mathcal{L}_{\mathrm{align}} \;=\; \mathrm{Wass}_1\!\bigl(\{g_{\bm w}\}_{\bm w\in
  \mathcal{S}}, \{g_{\bm w}\}_{\bm w \in \mathcal{T}_{\mathrm{adapt}}}\bigr),
\label{eq:align}
\end{equation}
added to the source cross-entropy with weight $\lambda$:
$\mathcal{L}=\mathcal{L}_{\mathrm{CE}}+\lambda\,\mathcal{L}_{\mathrm{align}}$.
Because the anchor is a scalar, the $1$-Wasserstein term needs only a sort;
it equates how the fault trajectory fills the phase space in the two domains and
is computationally free.

\subsection{C4: strict protocol and intervention audit}
\label{sec:c4}

Source and target never share a recording; the target is split into an unlabeled
\emph{adaptation} set (65\%) and a held-out \emph{test} set (35\%), also
recording-disjoint. We train on labeled source windows plus the target-adapt set
(for alignment) and report target-test accuracy. To evidence that the decision
\emph{uses} the geometry, we perturb descriptor $j$ at test time,
$\phi^{(j)}_e \leftarrow \phi^{(j)}_e + \delta\,\mathrm{std}_j\,z_e$ with
$z_e\sim\mathcal{N}(0,1)$, and measure the \emph{flip rate}, defined as the
fraction of test windows whose prediction changes. A decision that uses the geometry is
measurably disrupted by the perturbation, whereas an ignored descriptor leaves
decisions unchanged.

% Experimental setup.
\section{Experimental Setup}
\label{sec:setup}

\subsection{Data and cross-condition protocols}
\label{sec:data}

We use the CWRU bearing data \cite{zheng2019cwru}, which provides drive-end
acceleration under four motor loads (0--3\,hp) and ten classes (normal, inner-race,
ball, and outer-race faults at three severities each); load is the cross-domain
axis, shifting rotation speed and fault impulsiveness. Two disjoint protocols
carry the main evaluation:
\begin{itemize}
  \item \textbf{cwru\_03}: single-source hardest case. Train on CWRU load 0,
  transfer to load 3, ten classes.
  \item \textbf{cwru\_0123}: multi-source case. Train on loads 0,1,2, transfer
  to load 3, ten classes.
\end{itemize}
Both are recording- and condition-disjoint (Sec.~\ref{sec:method}), so no test
window shares a recording or operating condition with training.

\subsection{Models and implementation}
\label{sec:models}

We compare the raw-signal \textbf{1D-CNN} baseline \cite{ince2016real,
wen2018new,zhang2017new}, the source-only structural baseline
\textbf{DynHGNN+}, the full \textbf{GA-DCHG}, and its ablations, namely
GA-DCHG $-$align, $-$temporal, and $-$anchor. Embedding dimension $m=3$ chosen
by lag through mutual information \cite{fraser1986}, segments of $S=64$ state
points,
$n_{\mathrm{seg}}=4$ per window, neighbourhood order $k=8$, maximal hyperedge
size $K_{\max}=12$, anchor strength $\eta=1.0$, propagation $\tau=0.5$, and
alignment weight $\lambda=0.5$; we train for $60$ epochs with batch size $16$,
Adam with learning rate $2\times10^{-3}$, and the alignment term active from the
first epoch. Checkpoints are selected by source validation accuracy; every
protocol runs over three seeds $\{0,1,2\}$ and we report mean $\pm$ std
target-test accuracy.

% Results.
\section{Results}
\label{sec:results}

\subsection{R1: the geometric anchor is protocol-dependent}
\label{sec:r1}

Table~\ref{tab:screen} reports the normalized $1$-Wasserstein distance
$W_1(d)$ of each descriptor and the selected anchor for the two protocols.

\begin{table}[t]
  \centering
  \small
  \caption{Component 1: normalized source--target distance $W_1(d)$ of each
  edge-geometry descriptor and the data-selected geometric anchor.}
  \label{tab:screen}
  \begin{tabular}{lcc}
\toprule
Descriptor & cwru\_03 & cwru\_0123 \\
\midrule
$\kappa$ & 0.033 & 1.230 \\
$\chi$ & 0.070 & 1.097 \\
$\rho$ & 0.085 & 0.848 \\
\midrule
Anchor & $\kappa$ & $\rho$ \\
\bottomrule
\end{tabular}

\end{table}

The anchor is \emph{not} a fixed descriptor: curvature is the load-invariant
anchor for the single-source CWRU shift cwru\_03, where
$W_1=\WOneCZeroThreeCurvature$ is the smallest, whereas filling is the anchor
for the multi-source CWRU protocol cwru\_0123 with
$W_1=\WOneCZeroOneTwoThreeFilling$. Getting it wrong is concrete: on cwru\_03
filling is the \emph{most} domain-specific descriptor at
$W_1=\WOneCZeroThreeFilling$, so hard-coding filling as the anchor would inject
precisely the shift it is meant to remove; on cwru\_0123 curvature is the least
invariant at $W_1=\WOneCZeroOneTwoThreeCurvature$. The screening makes ``which
geometry is transferable'' measurable per protocol, using only unlabeled target
data.

\subsection{R2: geometric models are load-robust; GA-DCHG is best}
\label{sec:r2}

Table~\ref{tab:cross} reports target-test accuracy (mean $\pm$ std over three
seeds) and the gap to the 1D-CNN.

\begin{table*}[t]
  \centering
  \small
  \caption{Cross-domain target-test accuracy (mean $\pm$ std over three seeds)
  and the gap to the 1D-CNN baseline (positive = above the CNN).}
  \label{tab:cross}
  \begin{tabular}{lcccccc}
\toprule
Method & \multicolumn{3}{c}{cwru\_03} & \multicolumn{3}{c}{cwru\_0123} \\
\cmidrule(lr){2-4}\cmidrule(lr){5-7}
 & Acc. & std & Gap & Acc. & std & Gap \\
\midrule
1D-CNN & 0.560 & 0.126 & +0.000 & 0.788 & 0.109 & +0.000 \\
DynHGNN+ & 0.725 & 0.130 & +0.165 & 0.862 & 0.055 & +0.074 \\
\midrule
GA-DCHG & 0.753 & 0.125 & +0.193 & 0.902 & 0.031 & +0.114 \\
GA-DCHG $-align$ & 0.782 & 0.107 & +0.223 & 0.892 & 0.030 & +0.104 \\
GA-DCHG $-temporal$ & 0.768 & 0.107 & +0.209 & 0.872 & 0.057 & +0.084 \\
GA-DCHG $-anchor$ & 0.757 & 0.104 & +0.198 & 0.848 & 0.068 & +0.060 \\
\bottomrule
\end{tabular}

\end{table*}

\textbf{cwru\_03.} The raw 1D-CNN degrades sharply under the load shift, reaching
only \AccCZeroThreeCnn{} on the target test set; the structural baseline
DynHGNN+ is \GapCZeroThreeDynhg{} points \emph{above} the CNN, and GA-DCHG
improves on the structural baseline to reach the best accuracy,
\GapCZeroThreeGad{} points above the CNN (Table~\ref{tab:cross}). The
higher-order geometry of the reconstructed attractor is markedly more invariant
to the operating point than the raw waveform. The robustness that structural
models are assumed to lack is, on this protocol, precisely where they excel.

\textbf{cwru\_0123.} The multi-source protocol repeats the pattern: the 1D-CNN
reaches \AccCZeroOneTwoThreeCnn{}, DynHGNN+ is
\GapCZeroOneTwoThreeDynhg{} points above it, and GA-DCHG reaches
\AccCZeroOneTwoThreeGad{} (\GapCZeroOneTwoThreeGad{} above the CNN, the best of
all models). The consistent ordering across both protocols, GA-DCHG $\geq$
DynHGNN+ $\gg$ 1D-CNN, confirms the cross-domain advantage comes from the
phase-space geometry and that actively anchoring it helps.

\subsection{R3: component ablation}
\label{sec:r3}

Table~\ref{tab:ablation} isolates the three components on both protocols.

\begin{table*}[t]
  \centering
  \small
  \caption{Component ablation: target-test accuracy and the change ($\Delta$)
  relative to the full GA-DCHG.}
  \label{tab:ablation}
  \begin{tabular}{lcccc}
\toprule
Variant & \multicolumn{2}{c}{cwru\_03} & \multicolumn{2}{c}{cwru\_0123} \\
\cmidrule(lr){2-3}\cmidrule(lr){4-5}
 & Acc. & $\Delta$ & Acc. & $\Delta$ \\
\midrule
GA-DCHG & 0.753 & +0.000 & 0.902 & +0.000 \\
GA-DCHG $-align$ & 0.782 & +0.030 & 0.892 & -0.010 \\
GA-DCHG $-temporal$ & 0.768 & +0.015 & 0.872 & -0.030 \\
GA-DCHG $-anchor$ & 0.757 & +0.004 & 0.848 & -0.053 \\
\bottomrule
\end{tabular}

\end{table*}

On cwru\_0123 the component contributions are resolvable: removing the anchor
drops GA-DCHG by \XAccCZeroOneTwoThreeGad{}${}-{}$\XAccCZeroOneTwoThreeAnchor{}
to \XAccCZeroOneTwoThreeAnchor{}, and removing the temporal propagation to
\XAccCZeroOneTwoThreeTemp{}, at seed stds of only \XSdCZeroOneTwoThreeAnchor{} and
\XSdCZeroOneTwoThreeTemp{}. The anchor screening (C1) and temporal
propagation (C2) therefore carry part of the cross-domain benefit. The geometry-space
alignment (C3) is instead a zero-accuracy-cost, explainable regularizer:
removing it costs at most
\XAccCZeroOneTwoThreeGad{}$-{}$\XAccCZeroOneTwoThreeAlign{} points (cwru\_0123)
and nothing (cwru\_03), both within seed noise, and its value is the physical
grounding and auditability it gives the cross-domain coupling. On cwru\_03
the three ablations remain within seed-level noise of the full model (seed std up
to \SdMaxCZeroThree{}); at three seeds this single-source protocol alone cannot
separate them, but the multi-source protocol resolves the effect.

\subsection{R4: decisions are sensitive to the anchor geometry}
\label{sec:r4}

Table~\ref{tab:audit} reports the intervention audit: perturbing each descriptor
by $\delta\cdot\mathrm{std}$ at test time.

\begin{table}[t]
  \centering
  \small
  \caption{Intervention audit: flip rate when descriptor $j$ is perturbed at
  scale $\delta$ ($\ast$ marks the per-protocol anchor). Perturbing any geometry
  descriptor flips a non-trivial fraction of decisions, evidencing that geometry
  participates in the cross-domain decision.}
  \label{tab:audit}
  \begin{tabular}{llccc}
\toprule
Protocol & Descriptor & $\delta{=}0.5$ & $\delta{=}1.0$ & $\delta{=}2.0$ \\
\midrule
\multirow{3}{*}{cwru\_03} & $\kappa$$^{*}$ & 0.10 & 0.12 & 0.11 \\
 & $\chi$ & 0.12 & 0.11 & 0.11 \\
 & $\rho$ & 0.15 & 0.11 & 0.08 \\
\cmidrule(lr){1-5}
\multirow{3}{*}{cwru\_0123} & $\kappa$ & 0.10 & 0.09 & 0.10 \\
 & $\chi$ & 0.11 & 0.13 & 0.10 \\
 & $\rho$$^{*}$ & 0.11 & 0.09 & 0.10 \\
\bottomrule
\end{tabular}

\end{table}

Perturbing any of the three geometry descriptors flips a non-trivial fraction of
target-test decisions (\FlipScale{} on average), evidencing that the decision
depends on the geometry rather than only on the raw signal. The
descriptor-specific ranking is not cleanly separated at this sample size, so
the cross-domain decision is \emph{sensitive to the edge geometry} rather than
tied to one descriptor in particular, consistent with the screening's
identification of geometry as the transferable quantity.

% Discussion.
\section{Discussion}
\label{sec:discussion}

\textbf{Geometry is the transferable quantity.} The central result is that under
a load shift the higher-order geometry of the reconstructed phase-space attractor
transfers markedly better than the raw waveform: the structural models and
GA-DCHG sit \GapCZeroThreeDynhg{}/\GapCZeroThreeGad{} and
\GapCZeroOneTwoThreeDynhg{}/\GapCZeroOneTwoThreeGad{} points above the 1D-CNN on
cwru\_03 and cwru\_0123 respectively. Raw signal statistics such as amplitude
envelopes and energy bands are strongly tied to the operating point, whereas the local
curvature, expansion, and filling of the attractor are more intrinsic to the
fault's dynamical signature. GA-DCHG realizes this by making the geometry active,
and attains the best accuracy on both protocols.

\textbf{Anchor selection as method, not assumption.} Because the invariant
geometry is \emph{protocol-dependent}, hard-coding a descriptor inherits the
shift of a domain-specific one whenever the regime changes. Screening by the
normalized $1$-Wasserstein distance is simple, uses only unlabeled target data,
and is a standalone contribution that should transfer to other phase-space-based
representations.

\textbf{Intervention audit and the ablation trade-off.} R4 perturbs the edge
geometry at test time, rather than only measuring attention--geometry
correlation: perturbing the edge geometry flips
\FlipScale{} of decisions on average, showing the decision depends on the geometry.
The component ablation (R3) is within seed noise on
cwru\_03 but is resolvable on cwru\_0123, where removing the anchor or the
temporal propagation degrades the model beyond its seed std. The geometry-space
alignment (C3) acts as a zero-accuracy-cost, explainable regularizer: on
cwru\_0123 removing it costs at most 0.010 accuracy points, on cwru\_03 it
sits within seed noise, and its physical interpretation is independent of any
accuracy gain.

\textbf{Limitations.} We evaluate only the CWRU load axis, which has four loads,
at three seeds, bounding the resolution of component-level differences; a more
fine-grained multi-rig or multi-seed study would sharpen the component estimates.
We also note that an MFPT cross-load transfer is not feasible in the same form:
its normal class is recorded at a single load, and an extreme low\,$\rightarrow$\,high
fault-load shift transferred for no model, including the 1D-CNN, so we do not
report an MFPT headline number. Finally, the anchor is selected from unlabeled
target data, legitimate in unsupervised domain adaptation but to be revisited if
no target data were available.

\textbf{Broader impact.} Because the anchor is a scalar geometry statistic, the
pipeline is lightweight and runs on CPU for these data sizes; the screening and
intervention audit give practitioners a concrete way to decide \emph{which}
geometry transfers for their rig and to verify a deployed model actually uses it.
The scalar anchor is also relevant to uncertainty-aware, active model
updating \cite{hong2025bayes,song2025quad}, which can attach probabilistic
guarantees as operating conditions evolve.

% Conclusion.
\section{Conclusion}
\label{sec:conclusion}

We showed that the higher-order geometry of the reconstructed phase-space
attractor is precisely what transfers across a motor-load shift: on the CWRU
single- and multi-source protocols the structural models and GA-DCHG sit
\GapCZeroThreeDynhg{}/\GapCZeroThreeGad{} and
\GapCZeroOneTwoThreeDynhg{}/\GapCZeroOneTwoThreeGad{} points above a raw 1D-CNN,
which degrades sharply under the shift. Building on this, GA-DCHG uses the
edge-geometry descriptors as anchors rather than passive gates. The anchors
screen a per-protocol anchor, modulate hyperedge generation, align source and
target in geometry space, and validate decisions by intervention auditing.
GA-DCHG attains the best target-test accuracy of all models on both protocols.
The per-protocol anchor screening is itself a measurable result, the component
ablation on the multi-source protocol confirms the anchor and temporal components
carry part of the cross-domain benefit, and the intervention audit evidences
that the cross-domain decision is sensitive to the edge geometry. We release
all code so that every reported number is reproducible from the raw recordings.


\section*{Acknowledgements}
The authors thank the maintainers of the Case Western Reserve University
bearing data for making the raw recordings publicly available.

% References in Current Science style (numbered in order of citation).
% Surname, Initials., Title. Italic Journal, Year, Bold vol, pages.
\begin{thebibliography}{28}

\bibitem{li2021hg} Li, T., Zhou, Z., Li, S., Sun, C., Yan, R. and Chen, X., The emerging graph neural networks for intelligent fault diagnostics and prognostics: a guideline and a benchmark study. \textit{Mech. Syst. Signal Process.}, 2022, \textbf{168}, 108653.

\bibitem{li2021mrfcn} Li, T., Zhao, Z., Sun, C., Yan, R. and Chen, X., Multireceptive field graph convolutional networks for machine fault diagnosis. \textit{IEEE Trans. Ind. Electron.}, 2021, \textbf{68}, 12739-12749.

\bibitem{kipf2017} Kipf, T. N. and Welling, M., Semi-supervised classification with graph convolutional networks. In \textit{Proceedings of the International Conference on Learning Representations}, 2017.

\bibitem{feng2019} Feng, Y., You, H., Zhang, Z., Ji, R. and Gao, Y., Hypergraph neural networks. In \textit{Proceedings of the AAAI Conference on Artificial Intelligence}, 2019, \textbf{33}, 3558-3565.

\bibitem{ince2016real} Ince, T., Kiranyaz, S., Eren, L., Askar, M. and Gabbouj, M., Real-time motor fault detection by 1-D convolutional neural networks. \textit{IEEE Trans. Ind. Electron.}, 2016, \textbf{63}, 7067-7075.

\bibitem{wen2018new} Wen, L., Li, X., Gao, L. and Zhang, Y., A new convolutional neural network-based data-driven fault diagnosis method. \textit{IEEE Trans. Ind. Electron.}, 2018, \textbf{65}, 5990-5998.

\bibitem{zhang2017new} Zhang, W., Peng, G., Li, C., Chen, Y. and Zhang, Z., A new deep learning model for fault diagnosis with good anti-noise and domain adaptation ability on raw vibration signals. \textit{Sensors}, 2017, \textbf{17}, 425.

\bibitem{takens1981} Takens, F., Detecting strange attractors in turbulence. In \textit{Dynamical Systems and Turbulence, Lecture Notes in Mathematics}, Springer, Berlin, 1981, vol. 898, pp. 366-381.

\bibitem{packard1980geometry} Packard, N. H., Crutchfield, J. P., Farmer, J. D. and Shaw, R. S., Geometry from a time series. \textit{Phys. Rev. Lett.}, 1980, \textbf{45}, 712-716.

\bibitem{fraser1986} Fraser, A. M. and Swinney, H. L., Independent coordinates for strange attractors from mutual information. \textit{Phys. Rev. A}, 1986, \textbf{33}, 1134-1140.

\bibitem{kennel1992} Kennel, M. B., Brown, R. and Abarbanel, H. D. I., Determining embedding dimension for phase-space reconstruction using a geometrical construction. \textit{Phys. Rev. A}, 1992, \textbf{45}, 3403-3411.

\bibitem{grassberger1983} Grassberger, P. and Procaccia, I., Measuring the strangeness of strange attractors. \textit{Physica D}, 1983, \textbf{9}, 189-208.

\bibitem{ganin2016domain} Ganin, Y., Ustinova, E., Ajakan, H., Germain, P., Larochelle, H. et al., Domain-adversarial training of neural networks. \textit{J. Mach. Learn. Res.}, 2016, \textbf{17}, 1-35.

\bibitem{arjovsky2017wgan} Arjovsky, M., Chintala, S. and Bottou, L., Wasserstein generative adversarial networks. In \textit{Proceedings of the 34th International Conference on Machine Learning}, 2017.

\bibitem{gretton2012kernel} Gretton, A., Borgwardt, K. M., Rasch, M. J., Sch\"olkopf, B. and Smola, A., A kernel two-sample test. \textit{J. Mach. Learn. Res.}, 2012, \textbf{13}, 723-773.

\bibitem{wang2023hyp} Zhang, K., Li, H., Cao, S., Yang, C., Sun, F. and Wang, Z., Motor current signal analysis using hypergraph neural networks for fault diagnosis of electromechanical system. \textit{Measurement}, 2022, \textbf{201}, 111697.

\bibitem{zhao2020} Zhao, R., Yan, R., Chen, Z., Mao, K., Wang, P. and Gao, R. X., Deep learning and its applications to machine health monitoring. \textit{Mech. Syst. Signal Process.}, 2019, \textbf{115}, 213-237.

\bibitem{lei2020} Lei, Y., Yang, B., Jiang, X., Jia, F., Li, N. and Nandi, A. K., Applications of machine learning to machine fault diagnosis: a review and roadmap. \textit{Mech. Syst. Signal Process.}, 2020, \textbf{138}, 106587.

\bibitem{bao2025sthfd} Bao, P., Yi, W., Zhu, Y., Shen, Y. and Chai, B.-O., STHFD: spatial-temporal hypergraph-based model for aero-engine bearing fault diagnosis. \textit{Aerospace}, 2025, \textbf{12}, 612.

\bibitem{liu2026mat} Liu, D., Cheng, X., Jiang, Z. and Cui, L., Multi-attention collaborative temporal-spatial hypergraph attention network for machinery fault diagnosis. \textit{ISA Trans.}, 2026, \textbf{175}, 164-183.

\bibitem{li2025amh} Li, Y., Bwar, K. H., Chai, R., Tse, K. M. and Chai, B.-O., Adaptive multi-view hypergraph learning for cross-condition bearing fault diagnosis. \textit{Mach. Learn. Knowl. Extract.}, 2025, \textbf{7}, 147.

\bibitem{pan2010survey} Pan, S. J. and Yang, Q., A survey on transfer learning. \textit{IEEE Trans. Knowl. Data Eng.}, 2010, \textbf{22}, 1345-1359.

\bibitem{long2015learning} Long, M., Cao, Y., Wang, J. and Jordan, M. I., Learning transferable features with deep adaptation networks. In \textit{Proceedings of the 32nd International Conference on Machine Learning}, 2015, pp. 97-105.

\bibitem{zhao2025icmds} Zhao, Z., Li, J. and Cui, J., Feature extraction for gas turbine rotor fault diagnosis in complex networks based on phase space method. In \textit{Proceedings of the International Conference on Mechanical Design and Simulation (ICMDS)}, Springer, Xi'an, 2025.

\bibitem{vaswani2017} Vaswani, A., Shazeer, N., Parmar, N., Uszkoreit, J., Jones, L. et al., Attention is all you need. \textit{Adv. Neural Inf. Process. Syst.}, 2017, \textbf{30}.

\bibitem{zheng2019cwru} Smith, W. A. and Randall, R. B., Rolling element bearing diagnostics using the Case Western Reserve University data: a benchmark study. \textit{Mech. Syst. Signal Process.}, 2015, \textbf{64-65}, 100-131.

\bibitem{hong2025bayes} Hong, F., Wei, P., Bi, S. and Beer, M., Efficient variational Bayesian model updating by Bayesian active learning. \textit{Mech. Syst. Signal Process.}, 2025, \textbf{224}, 112113.

\bibitem{song2025quad} Song, J., Liang, Z., Wei, P. and Beer, M., Sampling-based adaptive Bayesian quadrature for probabilistic model updating. \textit{Comput. Methods Appl. Mech. Eng.}, 2025, \textbf{433}, 117467.

\end{thebibliography}


\end{document}
